We continue the source-cover analysis of linear-time transposition of a $d\times d$
matrix with $b(d)=\max(1,\lceil\log_2d\rceil)$-bit entries on a multitape Turing
machine. For a fixed admissible return geometry, the source-cover paper of this series
bounds the correctness deficit below by the mean circular gap area of the complete
raw-source banks, frozen at the first output time of each column, minus an explicit
revelation cost. We use this frozen gap area itself as the certificate. Along every
joint profile of every near-full sequence, its limiting mean over legal phases and
columns is at most $\beta\tau_R/(M-R)$, where $\tau_R$ is the limiting hop tail after
the first $R$ core families (Theorem~G). This implies the source-cover paper's
necessary reciprocal-cover bound. We apply Theorem~G to two classes of geometries, with
no envelope on the sizes of the core families.
A \emph{funnel} of depth $m$ is a chain of time sets in which each set is fed by at
most one family, and only from the next set, while the last set is fed by none. In no
near-full sequence do infinitely many geometries have funnels of a fixed depth that contain
every read of a raw interval of fixed positive relative length, whatever the fixed number of
exceptional families (Theorem~A). A \emph{single-sweep terminal-landing} geometry has two
designated families, an older and a newer one, and its exceptional families are among
them. It reads the input in one nondecreasing sweep during the lower region of the older
one, keeps the input head at one coordinate during the lower regions of the core families
and of the newer one and at every column root, and lets none of these families act during
the sweep. No near-full sequence consists of such geometries (Theorem~B). Theorem~G cannot
exclude the boundary case of the source-cover paper's power envelopes for two exceptional
families. An explicit sequence of admissible geometries with two exceptional families and
hop counts at most $512L/i^2$ has mean frozen gap areas at most $10^7/R^2$ at every
modulus and hop tails at least $1/(32R)$, so the necessary condition of Theorem~G holds on
it whenever $3.2\cdot10^8(M-R)\le\beta R$, for instance at $M=2R+2$ for every
$\beta\ge10^9$ (Theorem~Z). These geometries carry a single input: they are not near-full,
and they say what the certificate can see, not that a transposer exists. All numbered
theorems, lemmas, corollaries and definitions are formalized in Lean; the comparison with
the reciprocal-cover bound is a remark with a proof sketch only, and a few elementary
consequences stated in prose are labelled as not formalized. These are certificates for
specified geometries, not a decision of the transposition question.
